Este trabajo presenta un análisis experimental y comparativo de algoritmos fundamentales de ordenamiento (\textit{std::sort}, \textit{Merge Sort}, \textit{Quick Sort} y \textit{Patience Sort}) y multiplicación de matrices cuadradas (\textit{Naive} y \textit{Strassen}), contrastando sus cotas asintóticas teóricas ($\mathcal{O}, \Omega, \Theta, o, \omega$) frente a su rendimiento práctico de tiempo de CPU y memoria. Se evalúa la sensibilidad de los algoritmos ante el tamaño de las entradas, dominios y estructuras de datos. Los resultados demuestran que, en ordenamiento, las optimizaciones de bajo nivel y la localidad de memoria determinan la eficiencia empírica dentro de la clase $\Theta(n \log n)$, destacando a \textit{Patience Sort} por su fuerte dependencia del orden previo ($\Omega(n)$ vs $\mathcal{O}(n \log n)$). En multiplicación matricial, se evidencia cómo las constantes de operación y el acceso a la jerarquía de memoria favorecen al algoritmo iterativo tradicional en dimensiones moderadas, aproximándose al umbral de cruce teórico donde \textit{Strassen} ($o(n^3)$) comienza a amortizar su sobrecarga recursiva.